\chapter{Keratoderma}
\index{Keratoderma}

\section*{Definition and Overview}
Keratoderma is a group of skin conditions in which the skin on the palms of the hands and soles of the feet becomes thickened and hardened. The most common form is palmoplantar keratoderma, which may be inherited or acquired, and which can be diffuse, affecting the whole palm and sole, or focal, affecting pressure areas. The thickened skin can crack and become painful, affecting walking and hand function. Keratoderma can occur alone or as part of a syndrome with other skin or systemic features. Treatment focuses on moisturising, reducing thickening with keratolytic creams, and managing the underlying cause, with dermatology care for the more complex inherited forms.

\section*{Epidemiology}
Palmoplantar keratoderma is uncommon, and its inherited forms are rare, with a range of genetic conditions causing it. Acquired forms are more common and can follow other skin conditions, such as psoriasis and eczema, or systemic disease. It affects men and women similarly and can begin in childhood or adulthood depending on the cause. Because the palms and soles are constantly used, even mild keratoderma causes significant symptoms, and referral to a dermatologist is common.

\section*{Aetiology and Pathophysiology}
The skin normally produces new cells at the base and sheds them at the surface. In keratoderma, this process is disturbed so that the outermost layer (the stratum corneum) becomes excessively thick. In inherited forms, gene variants affect the production of keratin, the structural protein of skin, or the proteins that regulate skin turnover, causing the palms and soles to thicken. In acquired forms, chronic friction and pressure, or other skin diseases such as psoriasis, eczema, and lichen planus, cause thickening. Some systemic conditions, including hypothyroidism and certain cancers, can also cause acquired keratoderma. The thickened, inflexible skin cracks under pressure, causing pain.

\section*{Clinical Features}
\begin{itemize}
    \item Thickened, yellowed, or hardened skin on the palms and soles
    \item Diffuse involvement of the whole palm or sole, or focal calluses and cracks
    \item Cracking and fissuring, which can be deep and painful
    \item Pain with walking, standing, and gripping
    \item Onset in infancy or childhood in inherited forms
    \item Associated features: nail changes, thickened skin elsewhere, or syndromes in some forms
    \item Redness and scaling when associated with psoriasis or eczema
    \item Sweating abnormalities in some types
\end{itemize}

\section*{Differential Diagnosis}
Thickened palms and soles have several causes.
\begin{itemize}
    \item \textbf{Corns and calluses:} localised thickening from pressure, very common
    \item \textbf{Psoriasis:} palmoplantar psoriasis with scaling and redness
    \item \textbf{Eczema:} hand and foot dermatitis with itching
    \item \textbf{Fungal infection:} tinea of the palms and soles
    \item \textbf{Inherited keratodermas:} genetic forms, often with family history
    \item \textbf{Systemic conditions:} hypothyroidism, and rarely internal malignancy
    \item \textbf{Drug reactions:} thickening from some medicines
\end{itemize}

\section*{When to Seek Medical Care}
Thickening of the palms and soles that is progressive, painful, cracking, or affecting function should be assessed by a doctor, and dermatology referral is appropriate for diagnosis and treatment. Children with thickened palms or soles, or a family history, should be assessed early. Painful cracks that become red, swollen, or infected warrant prompt review.

\textbf{Contact your local emergency services for:}
\begin{itemize}
    \item Rapidly spreading redness, pain, and swelling, which may indicate cellulitis
    \item Fever with a skin infection
    \item Deep cracks that are bleeding heavily or very painful
    \item Signs of a severe systemic illness
\end{itemize}

\section*{Diagnosis and Investigations}
\begin{itemize}
    \item \textbf{History and examination:} the pattern, onset, family history, and associated skin disease
    \item \textbf{Skin biopsy:} to confirm the diagnosis and type
    \item \textbf{Genetic testing:} for inherited forms
    \item \textbf{Skin scrapings:} to exclude fungal infection
    \item \textbf{Blood tests:} for associated systemic conditions where indicated
    \item \textbf{Review of medicines and occupations:} contributing factors
    \item \textbf{Dermatology referral:} for specialised management
\end{itemize}

\section*{Management}
\begin{itemize}
    \item \textbf{Moisturising:} regular emollients to soften the skin
    \item \textbf{Keratolytics:} creams containing urea, salicylic acid, or lactic acid to reduce thickening
    \item \textbf{Retinoids:} oral medicines for severe inherited forms, under specialist care
    \item \textbf{Treatment of the underlying condition:} psoriasis, eczema, or fungal infection
    \item \textbf{Foot and hand care:} podiatry, cushioning, and appropriate footwear
    \item \textbf{Avoiding friction:} reducing pressure and using protective gloves
    \item \textbf{Treatment of cracks:} antiseptic care and dressings
    \item \textbf{Genetic counselling:} for inherited forms
\end{itemize}

\section*{Complications}
\begin{itemize}
    \item Painful fissures that interfere with walking and hand use
    \item Secondary bacterial infection of cracked skin
    \item Functional disability from thickened soles
    \item Cosmetic concerns and psychological effects
    \item Progression in inherited forms
    \item Complications of treatment, including irritation and systemic effects
\end{itemize}

\section*{Prognosis and Outcomes}
The outlook depends on the cause. Acquired keratoderma usually improves with treatment of the underlying condition and skin care. Inherited forms are lifelong, but most respond to regular moisturising, keratolytics, and, in severe cases, oral retinoids, allowing normal function. Dermatology care provides the best results, and most people with keratoderma manage their symptoms well with consistent treatment.

\section*{Prevention and Self-Care}
\begin{itemize}
    \item Moisturise the hands and feet daily, especially after washing
    \item Use keratolytic creams as directed to control thickening
    \item Wear cushioned footwear and well-fitted gloves
    \item Treat fungal infections and other skin conditions promptly
    \item Avoid harsh soaps and excessive friction
    \item Care for cracks to prevent infection
    \item Attend dermatology reviews for inherited forms
    \item Seek early treatment for new or worsening thickening
    \item Ask about genetic counselling if there is a family history
\end{itemize}

\section*{Sources and Further Reading}
The following reputable sources provide further information for healthcare students and patients seeking reliable, current advice about keratoderma.
\begin{itemize}
    \item Australasian College of Dermatologists. \textit{Skin condition information.} https://www.dermcoll.edu.au/
    \item Healthdirect Australia. \textit{Skin problems.} https://www.healthdirect.gov.au/
    \item DermNet NZ. \textit{Palmoplantar keratoderma.} https://dermnetnz.org/
    \item British Association of Dermatologists. \textit{Keratoderma patient information.} https://www.bad.org.uk/
    \item NHS. \textit{Thickened skin on hands and feet.} https://www.nhs.uk/
    \item National Organization for Rare Disorders. \textit{Palmoplantar keratoderma.} https://rarediseases.org/
\end{itemize}
